\documentclass[11pt]{article}
\usepackage[T1]{fontenc}
\usepackage{lmodern,microtype}
\usepackage[margin=1in]{geometry}
\usepackage{amsmath,amssymb,amsthm,mathtools}
\usepackage{booktabs,tabularx,longtable,array}
\usepackage{enumitem,xurl}
\usepackage[hidelinks]{hyperref}
\usepackage{fancyhdr}
\usepackage{listings}
\setlength{\headheight}{14pt}
\pagestyle{fancy}\fancyhf{}\fancyhead[L]{\small Certified disk-frontier reduction}\fancyhead[R]{\small ProveIt research}\fancyfoot[C]{\thepage}
\setlength{\parskip}{4pt}\setlength{\parindent}{0pt}
\setlist{nosep,leftmargin=1.6em}
\lstset{basicstyle=\ttfamily\small,breaklines=true,columns=fullflexible,keepspaces=true,frame=single,framerule=.3pt}
\newtheorem{theorem}{Theorem}[section]
\newtheorem{lemma}[theorem]{Lemma}
\newtheorem{proposition}[theorem]{Proposition}
\newtheorem{corollary}[theorem]{Corollary}
\theoremstyle{definition}\newtheorem{definition}[theorem]{Definition}
\theoremstyle{remark}\newtheorem{remark}[theorem]{Remark}
\newcommand{\F}{{\mathbb F_2}}
\newcommand{\Hom}{\operatorname{Hom}}
\newcommand{\im}{\operatorname{im}}
\newcommand{\id}{\operatorname{id}}
\newcommand{\rank}{\operatorname{rank}}
\newcommand{\End}{\operatorname{End}}
\newcommand{\Kh}{\operatorname{Kh}}
\newcommand{\code}[1]{\texttt{\detokenize{#1}}}
% BEGIN embedded measurements.tex
% Generated from final result files; do not edit by hand.
\newcommand{\DenseSpeedup}{5.87}
\newcommand{\SparseSlowdown}{8.79}
\newcommand{\UnitTests}{52}
\newcommand{\GeometryAccepted}{71}
\newcommand{\GeometryDeclined}{9}
\newcommand{\SmoothingChecks}{14,658}
\newcommand{\BenchmarkRepeats}{7}
\newcommand{\PythonVersion}{3.13.5}
\newcommand{\DenseBaseCompositions}{22,960}
\newcommand{\DenseNewCompositions}{5,658}

% END embedded measurements.tex

\title{Certified Disk-Frontier Reduction\ for Exact Unknot Recognition\\[5pt]
\large Common-boundary certificates, finite block cancellation,\ and two independent complexity obstructions}
\author{Research contribution prepared for Vladimir Reshetnikov's ProveIt project}
\date{7 October 2026}
\begin{document}
\maketitle
\begin{abstract}
The current ProveIt recognizer already has scalar-survivor prediction and
interruptible cancellation, but its maintained report identifies a missing
hypothesis for stronger radical transfer: partial-diagram matchings need not
share a disk boundary. We supply a constructive geometric bridge. A
rotation-system certificate recognizes processed regions whose frontier lies on
one boundary component and certifies a common cyclic order for every smoothing,
without enumerating those smoothings. Connected and co-connected prefixes pass;
a polynomial ear-insertion construction supplies such orders for every
biconnected plane crossing graph.

On certified stages we implement exact finite homological perturbation,
retaining nonzero maps between adjacent surviving degrees. The article proves
the geometric claims, justifies the required algebra with its provenance, and
provides typed chain-homotopy certificates, guarded adapters, an independent
crossing-cube oracle, and reproducible measurements. Dense synthetic blocks
improve by up to \DenseSpeedup{} times in the final paired run. Sparse examples
regress, and the diagram timings do not establish a recognition speedup.

Two explicit families separate the outstanding difficulties. Long-trefoil
products have constant frontier width but exponentially many explicit minimal
summands. Descending unknot diagrams with $n=m^2$ crossings force frontier width
at least $m/3$ in every crossing order, despite admitting disk-certified orders.
The resulting certified-scan bound is $\operatorname{poly}(n,1+R)2^{O(B)}$,
where $B$ is maximum boundary half-width and $R$ maximum minimal multiplicity.
No universal bound placing these parameters in the quasi-polynomial regime is
claimed. The continuation addresses a concrete applicability gap, not a general
solution of the asymptotic recognition problem.
\end{abstract}
\clearpage
{\small\setlength{\parskip}{0pt}\tableofcontents}
\newpage

\section{Scope, contribution, and relationship to the repository}
\label{sec:scope}
The objective remains an exact recognizer with a proved $n^{O(\log n)}$ bound.
The present contribution addresses a specific implementation and proof gap on the
Khovanov-scanning route. It is not a replacement for the separate hierarchy-based
program. Lackenby's July 2026 preprint analyzes a hierarchy process but explicitly
separates its step count from an adequately specified running-time analysis;
Section~9 also leaves its hierarchy parameters unbounded in that discussion
\cite{Lackenby2026}. Our argument does not fill those topological algorithmic gaps.

The live repository changed during this investigation. The audited version of
\path{synthesis/radical.tex} already contains survivor prediction, a degree-gap
shortcut, adaptive pausing, and a discussion of sharp radical nilpotence and
finite transfer \cite{RepoRadical}. Those algebraic results are \emph{background
for this continuation}, not discoveries claimed anew here. In particular,
canonical survivor counts and the sharp exponent below should not be counted as
additional independent breakthroughs when this package is integrated.

The contribution is the following bridge: derive a common-disk certificate from
the actual processed diagram, construct eligible orders for biconnected crossing
graphs, and use the certificate to justify full transfer on the current
interned-list scanner. We also construct a descending-grid unknot family that
rules out a universal polylogarithmic-width argument on raw diagrams. The
underlying graph and homological tools are classical; no claim of literature
priority is made for their ingredients or for the abstract perturbation lemma.

\begin{center}
\begin{tabularx}{\textwidth}{@{}p{.30\textwidth}X@{}}
\toprule
Item & Status in this article\\\midrule
Common-disk compiler & Proved, implemented, and checked against all smoothings of
small certified prefixes.\\
Connected-cut order construction & Proved for biconnected plane crossing graphs;
polynomial construction, not an optimal-width algorithm.\\
Full finite transfer & Exact on a supplied, checked disk category; geometry-bound
scanner implementation with safe decline and scalar fallback.\\
Canonical profiles and sharp radical exponent & Existing project theory and
standard algebraic foundations, rederived to make the assumptions explicit.\\
Performance & Paired synthetic improvements and sparse regressions;
no demonstrated full-pipeline improvement.\\
General quasi-polynomial recognition & Still requires additional structural and
representation bounds.\\\bottomrule
\end{tabularx}
\end{center}

The optional \code{openai/math} catalogue was consulted. No result advertised
there is used as an axiom or as an unverified shortcut in the proofs. The actual
foundations used here are the tangle-category formulation, local cancellation,
and homological perturbation \cite{Khovanov2002,BarNatan2007,Crainic2004}.

\section{Input, geometry, and coefficient conventions}
\label{sec:model}
A nonempty planar-diagram input is a list of four-tuples of edge labels, one tuple
per crossing, in a consistent cyclic order. Every edge label occurs exactly
twice. Opposite slots are the two strands. The rotation system must describe a
connected genus-zero diagram. The final recognition wrapper additionally checks
that the straight-through strand system has one component. Empty diagrams require
an explicit convention for isolated circles and are deliberately rejected by that
wrapper rather than guessed to mean an unknot.

Write $G$ for the plane multigraph of crossings and edges. Loops and parallel
edges are allowed. Let $P$ be a set of already processed crossings. A frontier
edge has exactly one endpoint in $P$. Its label appears once in the processed
list. If there are $2b$ such labels, $b$ is the \emph{boundary half-width}.
This is the complete frontier, not merely the four slots touched by the latest
crossing.

The scanner uses the Frobenius algebra
\[
 A=\F[x]/(x^2),\qquad
 \Delta(1)=1\otimes x+x\otimes1,\quad
 \Delta(x)=x\otimes x,\quad \epsilon(1)=0,\quad\epsilon(x)=1.
\]
After delooping, an object is a perfect matching of the frontier, with an integer
homological degree and a multiplicity. Quantum shifts are forgotten throughout
this article, as in the rank-only backend under study. Raw homological degrees
are not normalized by the negative-crossing shift; total rank is unaffected.

For two matchings $a,c$ on a common disk boundary, let $k(a,c)$ be the number of
circles in their overlay. A basis of $\Hom(a,c)$ is indexed by subsets of those
circles, specifying dots. The Python encoding places the coefficient of the
monomial with mask $S$ at integer bit position $S$. Thus a coefficient requires
up to $2^{k(a,c)}$ bits; it is not a machine-word scalar when $b$ grows. The
current \code{Planar} circle ordering is part of the interface. No coefficient
is transferred into a different circle-numbering convention.

A critical typing rule is that bit zero is a scalar identity coefficient only
when source and target matchings are equal. For different matchings the undotted
basis cobordism is generally a saddle-type map, not an identity. Every binary
projection and every identity shortcut in the implementation checks the matching
types as well as the coefficient bits.

\section{A common-disk certificate from boundary walks}
\label{sec:disk}
\subsection{The permutation certificate}
Replace each processed crossing by a small oriented disk and each processed
internal edge by an untwisted band. Keep a marked stub for each frontier edge.
This is a ribbon neighborhood $N(P)$. It retracts to the processed crossing graph;
frontier stubs do not change its Euler characteristic.

Number the darts at crossing $v$ by $4v,4v+1,4v+2,4v+3$. Let $\sigma$ rotate
the darts at each crossing by one slot. Let $\alpha_P$ exchange the two darts of
an internal edge and fix a frontier dart. The cycles of
\begin{equation}\label{eq:boundary-permutation}
 \phi_P=\sigma\alpha_P
\end{equation}
are the boundary walks of $N(P)$. Fixing a frontier dart models the turn around
the end of its stub. Let $v_P,e_P,f_P$ be the number of crossing disks, internal
bands, and boundary walks. For a connected prefix,
\begin{equation}\label{eq:ribbon-euler}
 2-2g=v_P-e_P+f_P.
\end{equation}
A boundary walk is \emph{marked} if it contains a frontier dart.

\begin{definition}[Disk certificate]
A disk certificate consists of the processed tuples, or a cryptographic digest
binding them to the certificate, the boundary-walk cycles, and their extracted
frontier labels. The checker requires that the processed graph be connected,
that \eqref{eq:ribbon-euler} give $g=0$, and that at most one boundary walk be
marked. At a nonempty frontier exactly one is marked. Its labels, in cyclic
order, give $\theta_P$. Rotation and reversal of that order are immaterial.
\end{definition}

\begin{theorem}[Geometric certification]\label{thm:disk}
Let the full input be a connected plane crossing diagram. A processed prefix
with a disk certificate is contained in a disk $D_P$ whose boundary meets the
remaining diagram exactly at the frontier labels in order $\theta_P$.
Every smoothing of the processed crossings induces a noncrossing matching in
that same order after closed circles are delooped. Consequently the entire
formal prefix complex, and not just its currently stored matching types, is
an object of the corresponding disk tangle category.
\end{theorem}
\begin{proof}
The connected genus-zero ribbon surface $N(P)$ has one complementary disk in
the sphere for each of its boundary components. Retain the complementary disk
adjacent to the marked boundary as the outside region, and fill all its other
boundary components. The resulting surface is a disk with the asserted boundary.

There are no unprocessed crossings hidden in one of the filled, unmarked
complementary disks. Otherwise, the component of the unprocessed crossing graph
there would have to connect to the rest of the full connected graph. It cannot
leave that complementary disk without meeting $N(P)$, and any edge meeting
$N(P)$ would be a frontier edge on its boundary, contradicting that it is
unmarked. This is where the connected full-diagram hypothesis is used.

Smoothing replaces the local crossing by disjoint arcs in its crossing disk.
Together with the internal bands these are embedded arcs and circles in $D_P$.
Disjoint arcs in a disk have a noncrossing endpoint matching. Removing closed
circles by the Frobenius delooping isomorphism does not alter those endpoint
pairings. Saddles, linear combinations, and local cancellation are relative to
the same boundary. Hence the formal complex and its homotopy manipulations
belong to the certified disk category.
\end{proof}

At a closed frontier, one may choose any complementary boundary as the outside;
there are no labels to order. Its coefficient category is simply $\F$.

\begin{proposition}[Certificate cost]\label{prop:certcost}
After indexing the edge labels, one prefix certificate can be built and checked
in linear time in its number of darts. Certifying all prefixes independently
costs $O(n^2)$ indexed graph operations and polynomial bit complexity.
The certificate is independent of the number of smoothing states and matching
multiplicities.
\end{proposition}
\begin{proof}
Build edge occurrences, the internal-edge involution, and adjacency in a scan
through the darts. Connectivity, permutation cycles, and extraction of marked
labels each require one further linear traversal. Summing prefix lengths proves
the second assertion. Arbitrarily long integer labels can first be relabelled;
their input bit length is included in that preprocessing cost.
\end{proof}

\subsection{A second, algebraic check}
The adapter also checks every \emph{distinct live matching} against $\theta_P$.
Use the positions of its endpoints in the order. Scanning from left to right,
push the far endpoint when opening an arc and require a matching stack pop when
closing it. This is a linear noncrossing test. It is a useful consistency check
between the geometric prefix and the scanner arrays. Theorem~\ref{thm:disk},
not this finite test alone, supplies provenance for all possible smoothing states.

\subsection{Why the certificate cannot be omitted}
The three pairings
\[
 (01)(23),\qquad(02)(13),\qquad(03)(12)
\]
are not simultaneously noncrossing in any cyclic order of four points. In every
such order one of the three is the crossing pairing. A putative gluing of the
three abstract types can give an odd value for twice the genus in the planar
composition engine. Pairing data by themselves therefore do not define a common
arc algebra. The implementation declines this situation \emph{before} applying
nilpotence or transfer. A failed certificate is not a knot verdict and is not
silently treated as an empty or acyclic complex.

\section{Constructing certified crossing orders}
\label{sec:orders}
\begin{lemma}[Connected cuts give one marked boundary]\label{lem:bond}
If both $G[P]$ and $G[V(G)\setminus P]$ are connected, then $P$ passes the disk
certificate.
\end{lemma}
\begin{proof}
The ribbon neighborhood of the connected prefix is a connected genus-zero
surface. Every vertex of the connected complement lies in one component of its
complement in the sphere. Every frontier edge approaches the boundary of that
same component. The complementary components of a connected planar ribbon
surface are disks with distinct boundary components. Therefore every frontier
stub lies on a single boundary component.
\end{proof}

An $st$-order is an ordering beginning at $s$ and ending at $t$ in which every
other vertex has both an earlier and a later neighbor. Such an order has connected
prefixes and suffixes: repeatedly follow an earlier neighbor to $s$, or a later
neighbor to $t$. We give an elementary construction to make its algorithmic
obligations explicit rather than relying on a black-box ordering theorem.

\begin{theorem}[Polynomial certified orders]\label{thm:ear}
Every biconnected plane crossing graph with at least two vertices has an order
all of whose proper prefixes pass Theorem~\ref{thm:disk}. An explicit ear-insertion
algorithm constructs and verifies such an order in polynomial time, with an
$O(n^3)$ bound in the indexed graph model used here. The one-crossing case is
handled directly.
\end{theorem}
\begin{proof}
Choose adjacent distinct vertices $s,t$ and initialize the ordered set to
$[s,t]$. Suppose the current ordered set is $S$ and is not all of $V(G)$.
Take any connected component $K$ of $G-S$. It has at least two distinct neighbors
in $S$: a unique neighbor would be an articulation separating $K$ from the
nonempty remainder of $S$.

Choose attachments $a,z$ with $a$ earlier than $z$. Within $K$, find a simple
path connecting a neighbor of $a$ to a neighbor of $z$. It may be a single
vertex. Insert this path immediately after $a$, oriented from $a$ toward $z$.
Every inserted vertex has an earlier and a later neighbor along the path.
Existing vertices retain their earlier and later neighbors, since their
relative order has not changed. The endpoints remain $s$ and $t$.

At least one vertex is inserted each time. Connectivity searches, path searches,
order updates, and a direct articulation check are polynomial; the simple
implementation is bounded by $O(n^3)$ graph operations. At completion the order
is an $st$-order. The preceding observation and Lemma~\ref{lem:bond} certify
every proper prefix. The implementation independently recomputes the boundary
walk certificates instead of trusting only the construction invariant.
\end{proof}

This construction does \emph{not} minimize frontier width. Prefer a good supplied
order when it is certified, or certify stages of the existing order only when a
block operation is contemplated. Converting to an arbitrary ear order can make
width worse. Articulated crossing graphs are declined by this constructor;
ordinary scanning and the repository's established decomposition methods remain
available. No unsupported articulation-to-connected-sum equivalence is assumed.

\section{Descending-grid unknots: a width obstruction on actual unknots}
\label{sec:grid}
The preceding theorem supplies admissible geometry, not polylogarithmic geometry.
The following explicit family prevents conflating those conclusions.

\begin{theorem}[Unknots with unavoidable wide raw frontiers]\label{thm:grid}
For each integer $m\ge2$ there is an explicitly constructible classical unknot
diagram $D_m$ with $n=m^2$ crossings whose crossing graph is biconnected and contains
the $m\times m$ grid as a spanning subgraph. In every crossing order, some prefix
has at least $m/3$ frontier edges. Thus its maximum boundary half-width is at
least $m/6=\Omega(\sqrt n)$.
\end{theorem}
\begin{proof}
Draw $m$ horizontal and $m$ vertical strands crossing in a square array. There
are $m^2$ crossings and $4m$ endpoints on the boundary of the square. Initially
these strands are $2m$ open paths. In the exterior disk on the sphere, join two
adjacent remaining endpoints belonging to different open paths. This merges
those paths and removes the two joined endpoints. Until only one path remains,
such an adjacent pair exists: each open path has exactly two remaining endpoints,
so if all cyclically adjacent endpoints belonged to the same path there could
be only one path. Join the final two endpoints to close it.

Joining endpoints adjacent in the remaining cyclic order gives a noncrossing
exterior matching. No additional crossings or vertices were added. The shadow
has one straight-through component, is planar and four-valent, and still contains
every internal grid edge. The grid is biconnected: it is the union of unit square
cycles joined along edges, and the union of two biconnected graphs sharing two
vertices is biconnected. Adding the exterior edges preserves this property.

Traverse the shadow from a marked point and make the first visit to every
crossing the overpass. This is a descending diagram. Cutting at the marked point,
lift the traversed arc with strictly decreasing height. Each first visit is then
above the second visit at its double point. A height-monotone arc can be straightened
by level-preserving translations, with its closing arc kept outside. Closing
therefore gives an unknot. This argument also explains the finite descending
certificate supplied by the generator.

It remains to prove the width claim, which is independent of the overpasses.
Let $S$ contain between $m^2/3$ and $2m^2/3$ vertices of the grid. If at least
$m/3$ rows contain vertices of both $S$ and its complement, those rows contribute
at least $m/3$ distinct horizontal cut edges. Otherwise there is a completely
occupied row and a completely unoccupied row: mixed rows alone cannot account
for either of the two sets' required sizes. In each column, the vertical path
between those two rows crosses the cut. This gives at least $m$ cut edges.
In either case $|\partial S|\ge m/3$. Every ordering has a prefix of size
$\lfloor m^2/2\rfloor$, within the stated range. Its grid cut edges are also
frontier edges of the crossing graph. Dividing their number by two proves the
half-width bound.
\end{proof}

This is a limitation of an \emph{order-only proof strategy on the unsimplified
input}, not a lower bound for unknot recognition and not an assertion that our
scanner takes exponential time on this family. Descending structure itself is
a certificate of triviality. Reidemeister simplification, semantic recognition,
or a succinct coefficient representation may bypass the wide raw cut entirely.
In particular, the theorem must not be used to exclude a quasi-polynomial
recognizer with such mechanisms.

The bundle generates $D_m$ for $1\le m\le8$, supplies their descending traversals,
and validates plane geometry and certified orders. Independent crossing-cube
homology was computed only for $m\le3$; its total rank is two. Larger ranks are
not reported as computed. The proof above supplies their unknot status.

% BEGIN embedded grid_table.tex
\begin{table}[htbp]
\centering\small
\begin{tabular}{@{}rrrrl@{}}\toprule
$m$ & Crossings & Input-order width & Ear-order width & Independent cube\\\midrule
1 & 1 & 0 & 0 & rank two\\
2 & 4 & 4 & 4 & rank two\\
3 & 9 & 4 & 6 & rank two\\
4 & 16 & 8 & 12 & not computed\\
5 & 25 & 8 & 12 & not computed\\
6 & 36 & 10 & 16 & not computed\\
7 & 49 & 10 & 18 & not computed\\
8 & 64 & 12 & 20 & not computed\\
\bottomrule\end{tabular}
\caption{Explicit descending-grid unknots. Width counts frontier points, not half-width. Ear insertion supplies certification but can worsen width; it is not an optimizer. All listed descending traversals are verified.}
\label{tab:grid}
\end{table}

% END embedded grid_table.tex


\section{The certified arc algebra and its finite radical}
\label{sec:radical}
The results in this section are algebraic infrastructure already discussed in the
maintained project report \cite{RepoRadical}. We include their proofs to state
precisely where the common-disk hypothesis enters. Khovanov's arc-algebra
construction supplies the categorical setting \cite{Khovanov2002}.

Fix $2b$ points in a prescribed common cyclic order. For a dotted basis morphism
$x_S:a\to c$ define its intrinsic weight by
\begin{equation}\label{eq:weight}
 w(x_S)=b-k(a,c)+2|S|.
\end{equation}
Every weight lies between zero and $2b$. Weight zero occurs only for the undotted
identity of a matching: $k(a,c)=b$ implies $a=c$. A matching has at most $b$ overlay
circles, so $\dim_\F\Hom(a,c)\le2^b$.

\begin{lemma}[Additivity of intrinsic weight]\label{lem:weight}
Composition is homogeneous for \eqref{eq:weight}. The positive-weight subspaces
form a two-sided ideal $J$, and the quotient category has only scalar identities
of the matching objects.
\end{lemma}
\begin{proof}
Consider basis inputs with $u,v$ dots. If a composite component has positive
genus, its Frobenius evaluation over $\F$ vanishes. Otherwise let $r$ be the number
of genus-zero components of the glued surface. The input disks and the $b$ gluing
bands give
\[
 k(a,b')+k(b',c)-b=2r-k(a,c),
\]
where $b'$ denotes the intermediate matching, not the half-width. On a surviving
component with $q$ output circles, no input dot gives $q-1$ output dots and one
input dot gives $q$; two dots give zero. Thus each nonzero output has
$k(a,c)-r+u+v$ dots. Substitution gives additivity of \eqref{eq:weight}.
The weight-zero characterization then proves the ideal and quotient assertions.
\end{proof}

\begin{theorem}[Sharp characteristic-two nilpotence]\label{thm:nilpotence}
For $b\ge1$, in the full disk matching category,
\[
 J^{2b}=0,\qquad J^{2b-1}\ne0.
\]
For any subcategory of its matching objects the same upper bound holds. On
endomorphisms of one matching, the positive ideal has exponent $b+1$ instead.
\end{theorem}
\begin{proof}
The weight bound immediately gives $J^{2b+1}=0$. Suppose a product of $2b$
positive homogeneous factors is nonzero. All factors must have weight one, and
the output must have weight $2b$. Each weight-one factor is undotted. The output
must be the fully dotted basis element of an endomorphism space, since weight
$2b$ forces $k(a,c)=b$ and $a=c$.

But a composite of undotted surfaces cannot have that term over $\F$. Positive
genus evaluates to zero; an undotted closed sphere evaluates to zero; and each
surviving genus-zero component with boundary leaves one of its output circles
undotted in each term. Hence the fully dotted output is impossible.

For sharpness, label the boundary $0,\ldots,2b-1$ and, for $0\le r\le b-1$, let
$a_r$ have pairs
\[
 (0,2r+1),\ (1,2),\ (3,4),\ldots,(2r-1,2r),
\]
followed by adjacent pairs on the remaining labels. These are noncrossing.
Compose the $b-1$ undotted weight-one morphisms from $a_0$ through $a_{b-1}$,
a dot on the first arc of $a_{b-1}$, and the reverse sequence of $b-1$ undotted
morphisms. The surface is connected: the successive elementary saddles join
all the initial arc components. Counting its disks and gluing bands gives Euler
characteristic $2-b$ and it has $b$ output boundary circles, so its genus is zero.
It carries exactly one dot and therefore evaluates to $x_1\cdots x_b\ne0$ in
$\End(a_0)$. There are $2b-1$ positive factors.

Finally $\End(a)=\F[x_1,\ldots,x_b]/(x_i^2)$. A product of $b+1$ positive
monomials vanishes, whereas $x_1\cdots x_b$ survives.
\end{proof}

The sharpness witness is algebraic. It does not assert that a scanned knot
prefix forces the maximum transfer depth. The program verifies these witnesses
through $b=8$. Neither the category nor its exponent is assigned to an arbitrary
uncertified family of abstract pairings.

\section{Scalar residues, minimal multiplicity, and what they do not retain}
\label{sec:residue}
For a bounded cochain complex $C$, retain the constant identity coefficient only
between equal matchings:
\[
 (d_0)_{vu}=\begin{cases}
 [1]d_{vu}\,\id_a,&\text{if both objects have matching }a,\\
 0,&\text{otherwise.}
 \end{cases}
 \qquad \delta=d-d_0.
\]
Here $[1]$ means coefficient of the undotted basis vector. By
Lemma~\ref{lem:weight}, $d_0^2=0$ and $\delta$ is $J$-valued. Write
\begin{equation}\label{eq:beta}
 \beta^h_a(C)=\dim_\F H^h(C_a,d_0),\qquad
 r(C)=\sum_{a,h}\beta^h_a(C).
\end{equation}
The summands $C_a$ contain the copies of matching $a$ only.

\begin{theorem}[Canonical minimal multiplicities]\label{thm:minimal}
Every bounded matching-valued complex that is homotopy equivalent to $C$ and
has differential in $J$ contains exactly $\beta^h_a(C)$ copies of $a$ in degree
$h$. In particular, every complete scalar-unit cancellation of $C$ leaves these
counts, independently of pivot order.
\end{theorem}
\begin{proof}
Reduction modulo $J$ preserves sums, composition, and chain-homotopy identities.
It therefore carries a homotopy equivalence to a homotopy equivalence of binary
complexes for each matching. A complex with differential in $J$ becomes a complex
with zero differential, whose homology dimensions are its object counts.
These dimensions must equal \eqref{eq:beta}.

For the last assertion, a coefficient $1+u$ in an endomorphism ring is a unit.
Indeed, in the square-free algebra over $\F$, $u^2=0$ for every positive sum
$u$, so $(1+u)^{-1}=1+u$. Complete unit cancellation terminates after removing
pairs of objects. If the final differential had a nonzero scalar residue, there
would still be a unit pivot. Its differential therefore lies in $J$.
\end{proof}

This is a uniqueness statement for \emph{multiplicities}, not for the surviving
basis or differential. It explains why pivot optimization changes fill and
composition work but cannot lower the final number of explicit generators for
a fixed processed tangle. Changing the processed prefix by changing the scan
order is a different operation.

\begin{proposition}[Profiles are not sufficient continuation states]\label{prop:notstate}
Scalar-residue homology alone does not determine the surviving differential or
its behavior under subsequent gluing.
\end{proposition}
\begin{proof}
Over $R=\F[x]/(x^2)$, the two complexes
\[
 R\xrightarrow{x}R,\qquad R\xrightarrow{0}R
\]
have the same residue profile, one object in each of two consecutive degrees.
They are not homotopy equivalent. Closing the arc turns multiplication by $x$
into a binary map of rank one on $A$, giving total homology dimension two;
the zero map gives total dimension four.

In particular a positive-weight map can produce a scalar matrix entry after a
frontier changes and a circle is delooped. The operation of taking scalar
residues is not a replacement for the gluing functor.
\end{proof}

If occupied residue degrees are nonconsecutive, the differential in a minimal
model must vanish. At a closed frontier it also vanishes, since $J=0$. These
are precisely the kinds of shortcuts already integrated in the repository
\cite{RepoRadical}. Full transfer is needed when those shortcuts do not apply.

\section{Constructive finite transfer with nonzero survivors}
\label{sec:transfer}
\subsection{Extracting a scalar contraction}
For each matching and degree split the binary complex as
\[
 C^h_a=B^h_a\oplus H^h_a\oplus L^h_a,
 \qquad d_0:L^h_a\overset{\sim}{\longrightarrow}B^{h+1}_a.
\]
Column elimination selects independent image vectors and their lifts. Kernel
vectors obtained from dependent columns are used to complement boundaries.
This constructs inclusion $i$, projection $p$, and homotopy $h$, with
\begin{equation}\label{eq:sdr0}
 pi=1,\quad d_0h+hd_0=1+ip,\quad hi=ph=h^2=0,\quad d_0i=pd_0=0.
\end{equation}
All their coefficients are binary multiples of identities of equal matchings.
The implementation uses bit-packed binary vectors and can verify all these
identities independently. No structured low-rank block has to be supplied by
the caller: the contraction is extracted from the actual scalar residue.

\begin{theorem}[Disk-certified finite transfer]\label{thm:transfer}
Let $C$ be a valid bounded complex on a certified disk frontier and let
\eqref{eq:sdr0} be a scalar contraction. Then
\begin{align}
 I&=(1+h\delta)^{-1}i,&
 D&=p\delta(1+h\delta)^{-1}i,\label{eq:transfer}\\
 H'&=(1+h\delta)^{-1}h,&
 P&=p(1+\delta h)^{-1}\label{eq:transfermaps}
\end{align}
are finite expressions and give a strong deformation retract of $(C,d)$ onto
$(H,D)$. The differential $D$ lies in $J$, so this is a minimal complex with the
multiplicities \eqref{eq:beta}. For $b\ge1$,
\begin{equation}\label{eq:finiteD}
 D=\sum_{j=0}^{2b-2}p\delta(h\delta)^ji.
\end{equation}
When all positive entries preserve their matching, at most $b$ positive factors
are needed instead. For $b=0$, $D=0$.
\end{theorem}
\begin{proof}
The operators $h,i,p$ have intrinsic weight zero. Each factor of $\delta$ is in
$J$, so Theorem~\ref{thm:nilpotence} makes the inverse series finite. Thus the
small-perturbation hypothesis of the homological perturbation lemma holds, and
\eqref{eq:sdr0} is its special deformation-retract hypothesis
\cite{Crainic2004}. Its formulas give \eqref{eq:transfer}--\eqref{eq:transfermaps}.
Every term of $D$ contains a positive factor, hence $D\in J$. Theorem~\ref{thm:minimal}
then identifies its multiplicities. A term in \eqref{eq:finiteD} contains $j+1$
positive factors; $2b$ such factors vanish. The one-matching bound follows in
the same way from the last assertion of Theorem~\ref{thm:nilpotence}.
\end{proof}

For clarity, the certificate checker verifies the full identities
\begin{gather}\label{eq:allchecks}
 d^2=D^2=0,\qquad dI=ID,\qquad Pd=DP,\qquad PI=1,\\
 dH'+H'd=1+IP,\qquad H'I=PH'=(H')^2=0.\notag
\end{gather}
One can also see the special side conditions directly. Put
$A=\delta(1+h\delta)^{-1}$. Then
$I=i+hAi$, $P=p+pAh$, and $H'=h+hAh$. Substituting these expressions makes
$PI=1$, $H'I=PH'=(H')^2=0$ immediate from \eqref{eq:sdr0}. The standard identity
$d_0A+Ad_0+AipA=0$ gives $D^2=0$ after multiplication by $p$ and $i$.
The other two chain-map identities and the homotopy identity are the remaining
perturbation identities. The implementation checks them using full categorical
matrix composition, separately from its column-wise transfer loop.

\subsection{The streaming algorithm}
For each residual basis column $i_j$, compute:
\begin{lstlisting}
t = i_j
D_j = 0
while t is nonzero:
    y = delta * t
    D_j = D_j + p * y
    t = h * y
\end{lstlisting}
Addition is XOR. Multiplication by $p$ and $h$ uses only binary incidence data;
expensive composition occurs in the multiplication by $\delta$. The old radical
matrix is never overwritten with a full ring-valued Schur complement. Its original
sparse edges are reused at every iteration.

The implementation enforces the finite factor bound and checks that the next
term vanishes. A term beyond the proved bound raises an arithmetic error;
it is never silently truncated. It also checks the common cyclic order and
coefficient-basis bounds. The default return value does not retain all lifted
columns of $I$. Full $I,P,H'$ are computed only for an explicit certificate
request, at higher memory and checking cost.

\subsection{A concrete two-term interpretation}
Suppose the source and target bases split into a cancellable block and surviving
objects and the differential is
\[
 d=\begin{pmatrix}A&B\\ C&E\end{pmatrix},\qquad
 A=A_0+U,
\]
where $A_0$ is invertible over $\F$ and $U,B,C,E$ have the appropriate positive
parts after scalar splitting. The surviving map is the usual typed Schur
complement $E+CA^{-1}B$. The finite transfer evaluates it without updating every
entry of the eliminated block. Crucially, the surviving map need not vanish.
The dense benchmark family below has adjacent survivor degrees and nonzero
surviving maps, so its successful reduction is not explained by the earlier
degree-gap or zero-survivor shortcut.

\section{Cost, memory, and the quasi-polynomial target}
\label{sec:cost}
Let $N$ be the number of explicit objects before a block reduction, $r$ the
number of survivors, and $E_\delta$ the number of nonzero entries in the original
radical differential. Let $S_h,S_p$ be the numbers of nonzero scalar entries in
the extracted homotopy and projection. Write $A_b$ for one coefficient addition
cost and $M_b$ for one correctly typed cobordism composition cost, including
its bit arithmetic. Let $G(N)$ bound binary contraction construction and checking.
Dense Gaussian elimination gives a conservative $G(N)=O(N^3)$ scalar-field
operation bound; bit packing improves implementation constants and word counts.

\begin{proposition}[Column-transfer cost]\label{prop:cost}
Ignoring optional full-certificate output and charging geometry compilation
separately, finite transfer can be implemented in
\begin{equation}\label{eq:kernelcost}
 G(N)+O\!\left(N+E_\delta+
 \ell r\bigl(E_\delta M_b+(S_h+S_p)A_b\bigr)\right)
\end{equation}
operations in the stated mixed cost model, where $\ell\le2b$ includes a final
zero check, or $\ell\le b+1$ when all positive entries preserve matching.
For $b=0$ only binary reduction is needed.
\end{proposition}
\begin{proof}
Each residual column requires at most $\ell$ products by the original radical
matrix. A product scans at most $E_\delta$ edges. Applying the scalar maps scans
at most $S_h+S_p$ incidences and adds their coefficient values. Sum over the $r$
columns and include extraction and input traversal.
\end{proof}

For dense matrices $E_\delta,S_h,S_p=O(N^2)$, expensive composition work is
$O(brN^2M_b)$ instead of a generic $O(N^3M_b)$ ring-valued elimination bound.
This separates binary work from expensive coefficient work; it does not magically
remove the $G(N)$ term. It is useful when $r\ll N$ and the scalar cancellations
would otherwise produce substantial fill. On a sparse diagonal matrix, ordinary
cancellation is cheaper.

A coefficient occupies at most $2^b$ bits. Even a naive implementation gives
$M_b=\operatorname{poly}(b)2^{O(b)}$: there are at most $4^b$ monomial pairs,
and a conservative bound for evaluating each output as a bit vector is
$\operatorname{poly}(b)2^b$ bit operations. The resulting loose bound
$\operatorname{poly}(b)8^b$ is sufficient here. Faster Frobenius/subset-convolution
kernels from the previous project work can improve $M_b$ independently of this
transfer algorithm. No fixed-word assumption is being used.

With a bounded product cache of $K$ entries, the streaming implementation needs,
beyond its input and the geometry engine's own caches, a conservative auxiliary
bound
\[
 O\!\left(N^2+(N+r^2+K)2^b\right)
\]
bits, up to index and object overhead. The default local product cache is capped
at 4096 entries and may be disabled. This bound does not pretend that the upstream
geometry and monomial caches are free. Full certificate output can require
$O(N^2 2^b)$ coefficient storage and substantially more checking time.

\begin{theorem}[Parameterized certified-scan bound]\label{thm:qp}
Consider a fixed disk-certified crossing order for a diagram with $n$ crossings.
Let $B$ be its maximum boundary half-width and let $R$ be the largest explicit
minimal survivor count after any stage, including the initial object. With exact
coefficient arithmetic, an explicit scanner using the reductions above has a
bound of the form
\begin{equation}\label{eq:globalcost}
 \operatorname{poly}(n,1+R)\,2^{O(B)}.
\end{equation}
Equivalently, for $\Theta=B+\log_2(1+R)$ the bound is
$\operatorname{poly}(n)2^{O(\Theta)}$. In particular
$\Theta=O((\log n)^2)$ implies $n^{O(\log n)}$ time.
\end{theorem}
\begin{proof}
Attaching one crossing gives two smoothings, each creating at most two new closed
circles. After delooping, each old object produces at most eight new objects.
Thus the pre-reduction object count is at most $8R$. The number of differential
entries is at most quadratic in that count. Crossing transfers have bounded
local topology and $2^{O(B)}$ coefficient cost.

At a disk frontier the number of matching types is at most $4^B$, by encoding a
noncrossing matching as balanced parentheses. Pairing, composition-plan, and
monomial caches can therefore be charged by $2^{O(B)}$ factors as well, rather
than assumed absent. Apply Proposition~\ref{prop:cost} at each stage, include
polynomial preprocessing and the certificate costs of Proposition~\ref{prop:certcost},
and sum over $n$ stages. Taking logarithms of the parameter factors proves the
last formulation.
\end{proof}

The theorem applies to full explicit multiplicities, not to an unspecified
compressed substitute for $R$. The existence theorem for certified orders does
not bound $R$ or force $B$ into the stated regime. Theorem~\ref{thm:grid} already
precludes a universal polylogarithmic raw width bound, even on unknots. The next
section shows why a small width alone is insufficient in the other direction.

\section{Constant-width exponential multiplicity, with an explicit local model}
\label{sec:trefoil}
The following example supports the already identified multiplicity obstruction
in the project report \cite{RepoRadical}. Its detailed local matrices make clear
which representation is obstructed.

Cut an edge of the positive three-crossing two-strand trefoil closure to obtain
a long trefoil. After the first two crossing reductions, the last crossing gives
a nine-object complex over $R=\F[x]/(x^2)$. In degrees zero through three its
ranks are $2,2,3,2$, and its matrices, with columns as sources, are
\begin{equation}\label{eq:trefoilmatrices}
 d^0=\begin{pmatrix}1&x\\1&x\end{pmatrix},\qquad
 d^1=\begin{pmatrix}x&x\\1&1\\0&0\end{pmatrix},\qquad
 d^2=\begin{pmatrix}x&0&x\\1&x&1\end{pmatrix}.
\end{equation}
The input PD, all earlier prefix data, and strong deformation-retract certificates
are supplied in \code{results/chain_certificates.json}; its verifier replays the
crossing transfers and separately checks the matrix identities.

Here is an explicit reduction of \eqref{eq:trefoilmatrices}. Name the bases
$(a,b)$, $(c,d)$, $(e,f,g)$, and $(h,i)$ respectively. Use new bases
\[
 (b+xa,a),\quad (c+d,c),\quad (xe+f,e+g,e),\quad (xh+i,h).
\]
They are invertible basis changes. The differential becomes the three identity
pairs
\[
 a\longmapsto c+d,\qquad c\longmapsto xe+f,\qquad e\longmapsto xh+i,
\]
with the other three generators $b+xa$, $e+g$, and $h$ closed. The surviving complex
is therefore three copies of one arc, in degrees $0,2,3$, with zero differential.

\begin{proposition}[Explicit fixed-width obstruction]\label{prop:trefoil}
Concatenating $k$ copies of this long trefoil gives a tangle with a zero-differential
minimal model containing $3^k$ copies of one arc. Its degree-count polynomial is
$(1+t^2+t^3)^k$. A sequential presentation uses at most six frontier points.
No fully explicit matching-valued homotopy-minimal representation can have fewer
than $3^k$ objects.
\end{proposition}
\begin{proof}
Concatenation is a planar composition functor and preserves each of the displayed
homotopy equivalences. Composing $k$ direct sums of three identity arcs multiplies
their multiplicities and their degree polynomials. The differential is zero.
The six-point bound follows by scanning one three-crossing long factor at a time,
retaining the outer endpoint while traversing the next factor; it is also checked
on the supplied diagrams. The final statement follows from
Theorem~\ref{thm:minimal}.
\end{proof}

This is not a lower bound for decision-only recognition. Closing the arc gives
rank $2\cdot3^k$, and structural recognition, component sharing, degree-polynomial
weights, or rank saturation can reject such inputs without materializing all
copies. Indeed, these are exactly the kinds of mechanisms already available in
the repository. The lesson is that matching type counts, coefficient dimensions,
and explicit multiplicities must be controlled separately.

\section{A guarded bridge to the current adaptive scanner}
\label{sec:integration}
The inspected \code{FastScan.eliminate} API accepts an optional
\code{update_budget} and returns whether cancellation completed
\cite{RepoScanner}. Its allowance counts candidate Schur update pairs, not
elapsed time and not necessarily expensive compositions. Before a pivot it
knows the Markowitz quantity
\[
 (|\operatorname{out}(u)|-1)(|\operatorname{in}(v)|-1).
\]
When a pivot would exceed the allowance, it pauses before that pivot, retains
all previously completed cancellations and valid maps, and clears its candidate
queues. This is the appropriate seam for the present work; restarting from the
original unreduced stage would throw away useful progress.

A conservative integration sequence is: preserve ordinary sparse cancellation;
use the existing residue and degree-gap shortcut when applicable; otherwise
attempt a common-disk certificate for the actual processed prefix; and only
then consider full finite transfer. If the certificate fails or a transfer
policy declines the block, resume the same partially reduced complex. The
mathematical validity of a pivot does not depend on which implementation performed
the preceding pivots.

The supplied \code{RadicalScan} demonstrates four research controls and the
new adaptive seam. Its \code{adaptive} mode starts with allowance
$\max(256,4N)$, preserves completed sparse progress, and attempts a geometry-checked
full reduction only after that allowance is exceeded. This fixture-level
adapter is not the repository's production adaptive mode and should not be
presented as such. Production already has additional profile shortcuts and
integration constraints \cite{RepoRadical}.

The controls are deliberately retained in the raw data. \code{always} tries
full transfer at every eligible stage; \code{auto} uses an exploratory survivor
fraction threshold; \code{terminal} uses binary completion only at closure.
They are ablations, not recommended defaults. Every open-boundary full transfer
requires the prefix certificate. Closed binary completion is an existing
principle, not a novel algorithm attributed to this package.

\paragraph{Atomicity and resource semantics.}
The block routine operates on a compact snapshot and installs its output only
after the computation succeeds. During an adaptive attempt, the authoritative
state is the already partially cancelled complex. Geometry decline preserves
it. Deadline or object-limit exceptions remain resource exceptions. A violated
algebraic identity raises an error rather than silently producing a verdict.
An externally supplied \code{update_budget} is delegated to the upstream
pair-budget operation and is never reinterpreted as a perturbation-depth budget.

The current wrapper supports the standard interned-list scanner. It does not
update component-owner arrays or saturated weight structures, install a Rust
implementation, modify race semantics, or alter the public CLI. Those integrations
need their own invariants and regression tests. The source-derived fixture is
kept under \code{reference_upstream}, separate from the original modules and from
production code.

\section{Correctness evidence and reproducibility}
\label{sec:validation}
The final bundle contains \UnitTests{} passing standard-library unit tests. The
test count is a count of test methods, many of which run families of cases; it is
not a count of formally verified theorems. Important checks include full scalar
contraction identities, typed coefficient bounds, invalid disk orders, explicit
radical witnesses, nontrivial three-term complexes, and failure semantics.

The unit suite compares 100 seeded small braid closures against an independent
crossing-cube implementation. That oracle directly forms smoothing circles and
uses $m$ and $\Delta$ over $\F$; it does not call \code{Planar}, its composition
plans, or its transfer routines. Both ordinary scanning and the guarded block
scanner are compared to its homological ranks. An additional 40 braid cases
compare every one of eight intermediate survivor profiles, giving 320 prefix
comparisons. Arbitrary shuffled orders are allowed in those tests precisely
because uncertified block stages fall back to the ordinary scanner.

A separate geometry experiment uses 80 seeded two-, three-, or four-strand braid
presentations with at most eight crossings. The ear-order constructor accepts
\GeometryAccepted{} biconnected cases and declines \GeometryDeclined{} others.
For every accepted prefix, \emph{all} smoothing states are formed independently
and their matchings are checked against the one computed cyclic order. This
produces \SmoothingChecks{} smoothing/order checks. The same
\GeometryAccepted{} closed diagrams are compared against the independent cube.
The log records cases, words, orders, widths, and full boundary certificates;
it does not claim that repeated presentations represent distinct knots.

Seven complete strong deformation-retract certificates are exported for the
four prefixes of a figure-eight presentation and the three prefixes of the long
trefoil. The checker verifies the coefficient and homological types, the disk
certificate, and all of \eqref{eq:allchecks}. It also replays the source-derived
crossing transfers to bind the certified input matrices to their listed PD
prefixes. This is executable exact algebra, not a Lean-checked development or
an independently verified implementation of every topological primitive.

The descending-grid generator checks its planar rotation systems and its complete
descending traversals through $m=8$. Independent cube ranks are additionally
computed through $m=3$. None of these finite experiments replaces the general
proof of Theorem~\ref{thm:disk}, \ref{thm:ear}, or \ref{thm:grid}.

\begin{center}
\begin{tabularx}{\textwidth}{@{}p{.40\textwidth}X@{}}
\toprule
Artifact & Reproduction command from bundle root\\\midrule
Unit suite & \code{python -m unittest discover -s tests -v}\\
Geometric exhaustive checks & \code{python experiments/geometry_validation.py}\\
Exported chain certificates & \code{python experiments/verify_certificates.py}\\
Descending-grid family & \code{python experiments/grid_family.py}\\
Paired timings & \code{python experiments/benchmark.py --repeats 7}\\
Article & \code{cd article; sh build.sh}\\\bottomrule
\end{tabularx}
\end{center}

The original full repository was not cloned or tested in this runtime. The fixture
was transcribed from inspected upstream sources, retaining the executable algorithms
and shortening commentary. It is not byte-identical to those files. Pinned blob
identifiers, attribution, and a script for a subsequent normalized-AST comparison
are included. The production suite, production recognition pipeline, Rust code,
and races were not run. Their independently reported test counts must not be
added to this bundle's results.

\section{Paired performance results, including regressions}
\label{sec:benchmarks}
The final measurement run uses \BenchmarkRepeats{} trials per case on
CPython \PythonVersion{}, Linux x86-64. Each trial starts with fresh geometry and
composition caches. Method order is alternated or rotated between trials.
Medians below are elapsed milliseconds from \code{perf_counter_ns}; all individual
nanosecond observations are retained in JSON, with a CSV summary.

\subsection{Dense two-term blocks with nonzero surviving differentials}
The source and target each have $m+q$ objects of one matching with $b$ arcs.
The scalar top-left $m\times m$ block is invertible and all other scalar entries
vanish. The remaining coefficients are seeded radical elements. Thus exactly
$2q$ objects survive in two adjacent degrees. The resulting differential is
nonzero in every listed case, including the sparse control. These are valid
algebraic complexes because there are only two degrees; they are not asserted
to be realized by the measured knot prefixes.

Every block is additionally checked outside the timed region by expanding the
complex over the regular binary representation of its square-free ring. The
original matrix, ordinary cancellation output, and transfer output have identical
homology dimensions in that independent expansion.

Timing starts after construction of the input coefficient matrix, matching
geometry, and ordinary scanner adjacency lists. It includes the full scalar
contraction and transfer in the new method, and the full cancellation loop in
the baseline. This models a ready-to-reduce stage, not diagram parsing.

% BEGIN embedded matrix_table.tex
\begin{table}[htbp]
\centering\small
\begin{tabular}{@{}rrrlrrr@{}}\toprule
$b$ & $m$ & $q$ & Pattern & Scalar (ms) & Transfer (ms) & Ratio\\\midrule
3 & 12 & 1 & dense & 1.371 & 0.839 & 1.63\\
3 & 24 & 1 & dense & 8.453 & 2.456 & 3.44\\
3 & 40 & 1 & dense & 26.532 & 5.735 & 4.63\\
4 & 24 & 1 & dense & 33.318 & 9.551 & 3.49\\
4 & 40 & 1 & dense & 151.327 & 25.780 & 5.87\\
4 & 24 & 4 & dense & 57.036 & 48.987 & 1.16\\
3 & 24 & 1 & mixed scalar & 11.800 & 2.912 & 4.05\\
4 & 40 & 1 & sparse & 0.043 & 0.374 & 0.11\\
\bottomrule\end{tabular}
\caption{Final paired stage-kernel timings. Ratio is scalar time divided by transfer time; below one is a regression. Source and target each have $m+q$ objects, with $2q$ survivors.}
\label{tab:matrix}
\end{table}

% END embedded matrix_table.tex


The largest observed median ratio is \DenseSpeedup{} on a dense block. On that
case the baseline requests \DenseBaseCompositions{} cobordism compositions,
whereas transfer requests \DenseNewCompositions{}. These counters count algorithm
requests under their documented identity/cache shortcuts; they are not wall-time
units. The survivor map remains nonzero. Thus the improvement extends beyond
an acyclic-block or degree-gap special case.

The sparse diagonal control goes the other way: the new method is
\SparseSlowdown{} times slower. Ordinary cancellation has no Schur fill there,
whereas binary basis extraction still costs work. The smaller speedup when $q$
is larger is consistent with the $r$ dependence in \eqref{eq:kernelcost}.
These observations favor a guarded dense-block path, not unconditional replacement.

\subsection{Complete fixed-order diagram scanners}
These measurements include every crossing transfer, elimination stage, and final
rank extraction, but exclude PD construction, input validation, order selection,
structural/polynomial prefilters, and whole-process startup. They are
\emph{complete fixed-order scanner timings}, not full production recognition
timings. The diagrams are braid closures and may be links; the comparison is of
homological ranks, not unsupported knot verdicts.

% BEGIN embedded diagram_table.tex
\begin{table}[htbp]
\centering\small
\begin{tabular}{@{}lrrrr@{}}\toprule
Presentation & Scalar (ms) & Eager (ms) & Adaptive (ms) & Block calls\\\midrule
Trefoil & 0.401 & 0.645 & 0.346 & 0\\
Figure-eight & 0.487 & 0.867 & 0.483 & 0\\
Torus $(3,5)$ & 3.067 & 5.227 & 3.302 & 0\\
Alternating 3-braid, 12 & 12.053 & 29.050 & 11.614 & 0\\
Cancelling 3-braid, 14 & 2.762 & 5.448 & 3.032 & 0\\
Mixed 4-braid, 12 & 5.284 & 8.550 & 5.222 & 0\\
Positive 4-braid, 12 & 8.276 & 12.314 & 8.357 & 0\\
\bottomrule\end{tabular}
\caption{Complete fixed-order scanner timings, with the same source-derived fixture and fresh caches. The last column counts adaptive full block transfers, not scalar cancellations.}
\label{tab:diagram}
\end{table}

% END embedded diagram_table.tex


For all seven diagram cases, the progress-preserving adaptive control completes
ordinary cancellation before its allowance is exhausted. It therefore invokes
\emph{zero full block transfers}. Its small timing differences from the baseline
cannot be attributed to the new algebra. The eager block control is slower on
every case. The exploratory survivor-fraction and terminal-only controls are
also retained in the raw data, without selecting their most favorable trial.

Accordingly, this run establishes a synthetic dense-kernel improvement and a
working, guarded diagram implementation. It does not establish that dense
cancellable blocks of the favorable kind arise often enough on difficult knots
to improve overall recognition. The absence of a measured diagram-level block
gain is an explicit result of the experiment, not omitted evidence.

\section{Why the final rank test is sound}
\label{sec:detector}
For a validated one-component classical diagram, the wrapper declares the unknot
exactly when unreduced Khovanov homology over $\F$ has total rank two. We recall
the characteristic-two reduction to the usual unknot-detection theorem rather
than assuming that every link of rank two has been validated as a knot.

Let $X$ multiply by $x$ on the circle containing a marked point. Let $\partial$
be the sum, over all circle factors, of the operation that replaces $x$ by $1$
and sends $1$ to zero. Direct checking on the four multiplication basis tensors
and two comultiplication basis vectors shows that $\partial$ commutes with the
crossing-cube differential over $\F$. Also
\[
 X^2=\partial^2=0,\qquad \partial X+X\partial=1.
\]
Hence the chain complex splits as the images of the complementary idempotents
$X\partial$ and $\partial X$. Those images are $\im X$ and $\im\partial$,
which are isomorphic as complexes via $X$ and $\partial$. The first is a reduced
complex, up to grading shifts. Consequently unreduced total rank over $\F$ is
twice reduced total rank.

If the unreduced rank is two, reduced rank over $\F$ is one. Universal coefficients
bound reduced rational rank above by that rank, while the reduced Euler
characteristic at $q=1$ has magnitude one for a knot. Reduced rational rank is
therefore one, and the unknot-detection theorem of Kronheimer and Mrowka applies
\cite{KM2011}. Conversely the unknot has one reduced and two unreduced generators.
This argument does not make the compressed coefficient representation itself a
new knot invariant. Its purpose is to preserve the existing invariant exactly.

\section{Further research questions and concrete next steps}
\label{sec:future}
The following are research directions and implementation obligations, not claims
that every item is an unsolved problem in the wider literature. Each mathematical
question deserves an independent literature audit before a novelty claim.

\begin{enumerate}[label=\textbf{Q\arabic*.},leftmargin=2.8em,itemsep=7pt]
\item \textbf{Output-sensitive dynamic disk certificates.}
Can the boundary permutation be maintained by local split/join operations as
crossings are added, using $O(n\log n)$ maintenance work plus the time needed to
emit requested cyclic orders? An explicit list of every frontier can itself
have quadratic total size, so an unconditional linear-size certificate is the
wrong target. The checker should replay local operations rather than trust a
mutable topology cache.

\item \textbf{Width-controlled certified ordering.}
Given a good existing order, characterize when it can be converted to a
connected-cut order without a large width increase. Measure this on real PDs
before using ear insertion as a performance heuristic. The descending-grid
family rules out a universal polylogarithmic raw-width conclusion; a useful
result must be relative to an instance parameter or include diagram simplification.

\item \textbf{Several marked boundary components.}
Develop a faithful, explicitly encoded category for prefixes whose frontier
occupies several boundary components. Which additional embedding or winding
information is needed? Treating all endpoint pairings as if they lay on one
disk is already falsified by the four-point example. A multisurface extension
must prove its composition and nilpotence claims in its actual category.

\item \textbf{Finding favorable adjacent-survivor blocks in practice.}
The new kernel is useful on dense complexes whose residual differential is
nonzero. Instrument difficult knot scans to determine whether such blocks arise,
how large $r/N$ is, and what fraction of wall time their reduction consumes.
Conway, Kinoshita--Terasaka, determinant-one examples, and scrambled certified
unknots are more informative than another isolated dense random matrix.

\item \textbf{A genuinely competitive adaptive policy.}
The current Schur-pair allowance is a transparent heuristic. Binary prediction
and full transfer can cost much more than the work completed before a pause.
A constant-factor portfolio theorem would need cooperative time slicing with
all primitive costs charged, or another proved domination argument. It does
not follow from the current switching rule.

\item \textbf{Geometry-aware scalar preflight.}
Can the constant part of a crossing transfer be compiled as a linear functional
on the incoming coefficient encoding, without expanding all outgoing coefficients?
Proposition~\ref{prop:notstate} shows that simply forgetting positive terms is
wrong. An exact preflight must account for the dots consumed by caps and cups
and preserve the source and target matching types.

\item \textbf{Combining two different forms of compression.}
Couple the component/Frobenius composition kernels and prior twist-block models
to the certified transfer engine. Kernel complexity $M_b$, explicit survivor
multiplicity $R$, and presentation length are separate parameters. A combined
bound must charge the extraction and maintenance of every compressed structure,
not only its multiplication after it has been supplied.

\item \textbf{Compression of nonzero minimal complexes.}
The long-trefoil obstruction can be handled with repeated zero-differential
summands, but generic minimal complexes have positive maps. Identify classes
closed under gluing and transfer that admit short descriptions of both their
multiplicities and those maps. Bounded-rank matrices, tensor descriptions, and
repeated components are candidates only when their maintained-size bounds can
be proved.

\item \textbf{Certificates with lower checking overhead.}
Full matrices $I,P,H'$ are easy to audit but can erase the memory advantage of
streaming. Can one verify block contraction from a succinct elimination or
factorization trace whose size is close to the work actually performed? Bind
that trace to the prefix rotation system and coefficient convention, not just
to an unlabelled binary matrix.

\item \textbf{Formalizing the trusted boundary.}
A useful formal target is the rotation-system disk theorem, the typed weight
calculation, and the finite perturbation identities over a certified finite
category. The present executable checks are not such a formalization. Negative
cases, especially incompatible cyclic orders and post-gluing scalar creation,
should be included in the specification rather than left as prose warnings.

\item \textbf{The hierarchy route remains separate.}
For the original quasi-polynomial objective, continue work on encoded layered
structures, compressed cutting with provenance, and constructive bounds for the
hierarchy operations. A polynomial disk-certificate compiler for a Khovanov
scanner is not a substitute for those algorithms or for their bit-cost analysis.

\item \textbf{A representation-sensitive end-to-end theorem.}
The decisive target is a constructive statement that every relevant input can
be simplified or decomposed into pieces whose \emph{maintained} description
entropy is $O((\log n)^2)$, with polynomially or quasi-polynomially many charged
operations. Theorem~\ref{thm:qp} describes the explicit-state version of that
obligation. The two obstruction families explain why a proof based solely on
pivot choice or solely on raw frontier order cannot supply it.
\end{enumerate}

\section{Conclusion and integration decision}
A common-disk hypothesis is now a checkable property of the processed diagram,
not an assumption inferred from the number of frontier labels. For biconnected
crossing graphs, a polynomial construction supplies orders with that property
at every prefix. This makes full finite transfer available on a rigorously
identified class of stages, including stages with nonzero maps between adjacent
survivor degrees.

The measured benefit is narrower than the mathematical applicability: favorable
dense algebraic blocks improve, sparse blocks regress, and the measured ordinary
diagrams do not exercise the profitable adaptive branch. Integration should
therefore preserve the current default, put the geometry check at the existing
pause boundary, and keep the new path opt-in until a complete-checkout regression
and representative workload study have been performed.

The asymptotic target remains substantive. Certified disk geometry, bounded
coefficient cost, and bounded explicit multiplicity are different requirements.
The descending grids and long-trefoil products separate them on concrete inputs.
Progress toward a general quasi-polynomial recognizer must address maintained
representation size or topological simplification, not merely choose a different
order of elementary cancellations.

\appendix
\section{Package map and source provenance}
The package contains the article source and compiled PDF, standard-library
implementations under \code{src}, independent and adversarial tests under
\code{tests}, experiment drivers, JSON certificates, raw paired measurements,
and integration notes. The principal entry points are:
\begin{lstlisting}
src/disk_frontier.py     # rotation-system certificates and ear orders
src/radical.py           # binary contraction and finite typed transfer
src/radical_scan.py      # guarded current-API scanner adapter
src/certified_driver.py  # validated homology / knot wrapper
src/descending_grid.py   # explicit certified unknot family
\end{lstlisting}

The audited source-derived fixture is based on these Git blob identifiers:
\begin{center}
\begin{tabular}{@{}ll@{}}\toprule
Module & Blob identifier prefix\\\midrule
\code{scan_fast.py} & \code{6c731d0f4d2c2c8bbe51}\\
\code{planar.py} & \code{1078526e7e7dbaf0b726}\\
\code{geometry.py} import surface & \code{96acc160dbf2df7981b2}\\
Maintained \code{synthesis/radical.tex} & \code{53da077349498476001f}\\\bottomrule
\end{tabular}
\end{center}
Full identifiers and the scope of transcription are in \code{PROVENANCE.md}.
The fixture's geometry module intentionally contains only the constants and
exception imported by the two scanner modules; it is not the upstream geometry
module in full. No claim of a complete checkout or automatic source-equivalence
verification is made. The supplied fetch-and-AST audit script is an additional
integration check to run in an environment with repository network access.

\section{Proof and evidence ledger}
\begin{longtable}{@{}p{.29\textwidth}p{.65\textwidth}@{}}
\toprule
Claim & Justification and boundary\\\midrule\endhead
Disk-frontier applicability & Theorem~\ref{thm:disk}; connected full plane input,
connected genus-zero prefix, one marked boundary.\\
Certified orders & Theorem~\ref{thm:ear}; biconnected crossing graph;
no low-width guarantee.\\
Grid width lower bound & Theorem~\ref{thm:grid}; raw diagrams; does not constrain
simplification or decision certificates.\\
Minimal counts & Theorem~\ref{thm:minimal}; fixed certified category; basis and
differential are not claimed canonical.\\
Block equivalence & Theorem~\ref{thm:transfer}; valid $d^2=0$, typed geometry,
verified scalar contraction, no silent truncation.\\
Quasi-polynomial specialization & Theorem~\ref{thm:qp}; conditional on
$B+\log(1+R)=O((\log n)^2)$; no universal extraction theorem.\\
Dense speedup & Final seven-trial fixture benchmark; synthetic complexes,
nonzero surviving maps; not a production pipeline comparison.\\
Full scanner correctness evidence & Independent cube comparisons and exported
chain certificates; finite software validation, not formal verification.\\
General recognition speedup & Not established by the measurements.\\
Production, Rust, race integration & Not performed.\\\bottomrule
\end{longtable}

\clearpage
\begin{thebibliography}{9}
\bibitem{Khovanov2002}
M. Khovanov, \emph{A functor-valued invariant of tangles}, Algebraic \& Geometric
Topology \textbf{2} (2002), 665--741.
\url{https://arxiv.org/abs/math/0103190}.
\bibitem{BarNatan2007}
D. Bar-Natan, \emph{Fast Khovanov homology computations}, Journal of Knot Theory
and Its Ramifications \textbf{16} (2007), 243--255.
\url{https://arxiv.org/abs/math/0606318}.
\bibitem{Crainic2004}
M. Crainic, \emph{On the perturbation lemma, and deformations}, 2004,
especially Section~2.
\url{https://arxiv.org/abs/math/0403266}.
\bibitem{KM2011}
P. B. Kronheimer and T. S. Mrowka, \emph{Khovanov homology is an unknot-detector},
Publications Math\'ematiques de l'IH\'ES \textbf{113} (2011), 97--208.
\url{https://arxiv.org/abs/1005.4346}.
\bibitem{Lackenby2026}
M. Lackenby, \emph{Incompressible surfaces, hierarchies and unknot recognition},
arXiv:2607.23350v1, 25 July 2026, especially Section~9.
\url{https://arxiv.org/abs/2607.23350}.
\bibitem{RepoRadical}
V. Reshetnikov's ProveIt repository, maintained synthesis,
\emph{Survivor prediction and adaptive cancellation},
\code{Topology/UnknotRecognition/synthesis/radical.tex}, inspected 7 October 2026;
Git blob \code{53da077349498476001f6b45065b34f0e029daa6}.
\url{https://github.com/VladimirReshetnikov/ProveIt/blob/main/Topology/UnknotRecognition/synthesis/radical.tex}.
\bibitem{RepoScanner}
V. Reshetnikov's ProveIt repository,
\path{Topology/UnknotRecognition/fast/fastunknot/scan_fast.py}, inspected source
with interruptible cancellation; Git blob
\code{6c731d0f4d2c2c8bbe513e6b7495cf5dbdc8dc88}.
\url{https://github.com/VladimirReshetnikov/ProveIt/blob/main/Topology/UnknotRecognition/fast/fastunknot/scan_fast.py}.
\end{thebibliography}
\end{document}
